% ==================================================
\appendix
\section{Weryfikacja numeryczna}
\label{sec:verification}
% ==================================================


\section*{A.1. Warunki eksperymentu}
\label{subsec:a1-warunki}

Eksperymenty przeprowadzono w Python~3.14 (SymPy~1.14, NumPy)
na stacji AMD Ryzen~9~5900X (64\,GB, Windows~11~Pro)
i laptopie Intel Core i7-13620H (16\,GB, Windows~11~Pro).
Arytmetyka: IEEE~754 podwójna precyzja ($\varepsilon \approx 2{,}2 \times 10^{-16}$).

Trzy implementacje --- naiwna $O(N^2)$, sumy prefiksowe $O(N\log N)$
i przesuwne okno $O(N)$ --- tworzą łańcuch weryfikacji.
Liczby operacji są deterministycznymi funkcjami $(a, b)$.
Sprawdzenie krzyżowe przez bezpośrednią enumerację $O(2^n)$
przy $n = 27, 28$ przeprowadzono wielowątkowo;
wszystkie pozostałe wartości całkowite
($w(n)$, liczby operacji, pozycje aktywne)
były identyczne bitowo na obu platformach.
Wszystkie skrypty weryfikacyjne dostępne są w repozytorium
\url{https://github.com/Flaman55/RelMathApps/tree/main/number_theory/numerical_semigroups/linear_Q}:
\texttt{verify\_linear\_Q.py} (Twierdzenia~\ref{thm:linear}--\ref{thm:general},
Tabele~\ref{tab:two-gen}--\ref{tab:multi-gen}),
\texttt{appendix\_verify.py} (Tabele~\ref{tab:mertens}--\ref{tab:sequences}),
\texttt{incremental\_windows.py} (algorytm inkrementalny),
oraz \texttt{section\_C\_parallel.py} (obliczenia równoległe, $n=19,\ldots,30$).


\section*{A.2. Dwa generatory}
\label{subsec:a2-dwa}

\begin{table}[H]
\centering
\begin{tabular}{rrrrrr}
\toprule
$a$ & $b$ & $N$ & oper.\ $O(N^2)$ & oper.\ $O(N)$ & maks.\ błąd \\
\midrule
  3 &   5 &     4 &          16 &        16 & $5{,}3\times10^{-15}$ \\
  5 &   7 &    12 &         144 &        48 & $1{,}8\times10^{-14}$ \\
  7 &  11 &    31 &         961 &       124 & $1{,}1\times10^{-13}$ \\
 11 &  13 &    60 &        3600 &       240 & $2{,}8\times10^{-13}$ \\
 23 &  29 &   314 &       98596 &      1256 & $1{,}2\times10^{-11}$ \\
 37 &  41 &   721 &      519841 &      2884 & $4{,}3\times10^{-11}$ \\
 97 & 101 &  4801 &   23049601 &     19204 & $1{,}3\times10^{-9}$  \\
\bottomrule
\multicolumn{6}{l}{\footnotesize Wszystkie 450 par $(a,b)$ z $b \leq 40$: \textbf{PASS}.
  Tolerancja $10^{-9}$ ($10^{-8}$ dla $N > 500$).}
\end{tabular}
\caption{Weryfikacja Twierdzenia~\ref{thm:linear}: dwa generatory.
Stosunek operacji $N^2/(4N) = N/4$.}
\label{tab:two-gen}
\end{table}


\section*{A.3. Wiele generatorów}
\label{subsec:a3-wiele}

\begin{table}[H]
\centering
\begin{tabular}{lrrrr}
\toprule
Generatory & $N$ & $w(n)$ & maks.\ błąd & Status \\
\midrule
$\{2,3\}$     & 1 & 2 & $0$                   & OK \\
$\{3,5\}$     & 4 & 2 & $8{,}9\times10^{-16}$ & OK \\
$\{2,3,5\}$   & 1 & 4 & $0$                   & OK \\
$\{3,5,7\}$   & 3 & 5 & $1{,}8\times10^{-15}$ & OK \\
$\{2,3,5,7\}$ & 1 & 8 & $0$                   & OK \\
\bottomrule
\end{tabular}
\caption{Weryfikacja Twierdzenia~\ref{thm:general}: wiele generatorów
wobec naiwnego $O(N^2)$.}
\label{tab:multi-gen}
\end{table}

\newpage
\section*{A.4. Dane gęstości (Twierdzenie~\ref{thm:density})}
\label{subsec:a4-gestosc}

\begin{table}[H]
\centering
\small
\begin{tabular}{r|rrrrr}
\toprule
$n$ & $2^n-1$ & $w(n)$ & $P(n)$ & Rzadkość\% & $\delta(n)$ \\
\midrule
 2 &          3 &    2 & 0{,}3333 &  0{,}0\% & $-0{,}3333$ \\
 3 &          7 &    4 & 0{,}2667 & 14{,}3\% & $-0{,}3048$ \\
 4 &         15 &    8 & 0{,}2286 & 26{,}7\% & $-0{,}3048$ \\
 5 &         31 &   14 & 0{,}2078 & 29{,}0\% & $-0{,}2438$ \\
 6 &         63 &   20 & 0{,}1918 & 44{,}4\% & $-0{,}1257$ \\
 7 &        127 &   30 & 0{,}1805 & 59{,}1\% & $-0{,}0557$ \\
 8 &        255 &   38 & 0{,}1710 & 72{,}2\% & $+0{,}0220$ \\
 9 &        511 &   48 & 0{,}1636 & 81{,}6\% & $+0{,}0697$ \\
10 &       1023 &   60 & 0{,}1579 & 88{,}0\% & $+0{,}0993$ \\
11 &       2047 &   84 & 0{,}1529 & 92{,}5\% & $+0{,}1118$ \\
12 &       4095 &  116 & 0{,}1487 & 95{,}3\% & $+0{,}1204$ \\
13 &       8191 &  146 & 0{,}1451 & 97{,}2\% & $+0{,}1273$ \\
14 &      16383 &  198 & 0{,}1417 & 98{,}3\% & $+0{,}1296$ \\
15 &      32767 &  236 & 0{,}1387 & 99{,}0\% & $+0{,}1315$ \\
16 &      65535 &  288 & 0{,}1361 & 99{,}4\% & $\mathbf{+0{,}1317}$ \\
17 &     131071 &  352 & 0{,}1338 & 99{,}7\% & $+0{,}1311$ \\
18 &     262143 &  412 & 0{,}1316 & 99{,}8\% & $+0{,}1300$ \\
19 &     524287 &  492 & 0{,}1296 & 99{,}9\% & $+0{,}1287$ \\
20 &    1048575 &  564 & 0{,}1278 & 99{,}9\% & $+0{,}1273$ \\
21 &    2097151 &  646 & 0{,}1260 &100{,}0\% & $+0{,}1257$ \\
22 &    4194303 &  720 & 0{,}1245 &100{,}0\% & $+0{,}1243$ \\
23 &    8388607 &  808 & 0{,}1230 &100{,}0\% & $+0{,}1229$ \\
24 &   16777215 &  892 & 0{,}1216 &100{,}0\% & $+0{,}1215$ \\
25 &   33554431 &  996 & 0{,}1203 &100{,}0\% & $+0{,}1203$ \\
26 &   67108863 & 1096 & 0{,}1191 &100{,}0\% & $+0{,}1191$ \\
27 &  134217727 & 1200 & 0{,}1180 &100{,}0\% & $+0{,}1180$ \\
28 &  268435455 & 1306 & 0{,}1169 &100{,}0\% & $+0{,}1169$ \\
29 &  536870911 & 1416 & 0{,}1158 &100{,}0\% & $+0{,}1158$ \\
30 & 1073741823 & 1528 & 0{,}1148 &100{,}0\% & $+0{,}1148$ \\
\bottomrule
\end{tabular}
\caption{Liczba okien aktywnych $w(n)$, iloczyn gęstości $P(n)$,
rzadkość $1 - w(n)/\sigma_n$
i deficyt $\delta(n)$ (Twierdzenie~\ref{thm:density})
dla $p_i = i$-ta liczba pierwsza.
Wartości graniczne: $n=8$ --- pierwsze $\delta(n) > 0$;
$n=21$ --- rzadkość osiąga $100\%$ (po zaokrągleniu);
$n=25$ --- $w(n)/(2^n-1) < 10^{-4}$.
Maksimum $\delta(n)$ przy $n=16$ (pogrubione).
Wartości ,,100,0\%'' wynikają z zaokrąglenia do jednego miejsca dziesiętnego.
Dane $n \leq 18$: enumeracja bezpośrednia;
$n = 19, \ldots, 30$: algorytm inkrementalny,
niezależnie potwierdzony dla $n=27,28$ enumeracją $O(2^n)$.}
\label{tab:mertens}
\end{table}

\newpage
\section*{A.5. Wartości $K^{(n)}(d)$ --- aktywacja i dezaktywacja}
\label{subsec:a5-wartosci}

\begin{table}[H]
\centering
\small
\begin{tabular}{r|rrrrrrrrrrr}
\toprule
$d$ & $n{=}2$ & $n{=}3$ & $n{=}4$ & $n{=}5$ & $n{=}6$ &
      $n{=}7$ & $n{=}8$ & $n{=}9$ & $n{=}10$ & $n{=}11$ & $n{=}12$ \\
\midrule
 0 &  1 &  1 &  1 &  1 &  1 &  1 &  1 &  1 &  1 &  1 &  1 \\
 3 & -1 & -1 & -1 & -1 & -1 & -1 & -1 & -1 & -1 & -1 & -1 \\
 5 &  $\cdot$ & -1 & -1 & -1 & -1 & -1 & -1 & -1 & -1 & -1 & -1 \\
 7 &  $\cdot$ &  $\cdot$ & -1 & -1 & -1 & -1 & -1 & -1 & -1 & -1 & -1 \\
 8 &  $\cdot$ &  1 &  $\cdot$ &  $\cdot$ &  $\cdot$ &  $\cdot$ &  $\cdot$ &  $\cdot$ &  $\cdot$ &  $\cdot$ &  $\cdot$ \\
 9 &  $\cdot$ &  $\cdot$ &  1 &  1 &  1 &  1 &  1 &  1 &  1 &  1 &  1 \\
10 &  $\cdot$ &  $\cdot$ &  1 &  1 &  1 &  1 &  1 &  1 &  1 &  1 &  1 \\
16 &  $\cdot$ &  $\cdot$ &  $\cdot$ &  $\cdot$ &  1 &  1 &  1 &  1 &  1 &  1 &  1 \\
17 &  $\cdot$ &  $\cdot$ &  $\cdot$ &  1 &  2 &  1 &  1 &  1 &  1 &  1 &  1 \\
25 &  $\cdot$ &  $\cdot$ &  $\cdot$ &  $\cdot$ &  $\cdot$ &  $\cdot$ &  1 &  1 &  1 &  1 &  1 \\
47 &  $\cdot$ &  $\cdot$ &  $\cdot$ &  $\cdot$ &  $\cdot$ &  1 &  1 &  $\cdot$ & -1 & -2 & -3 \\
\bottomrule
\end{tabular}
\caption{Wybrane wartości $K^{(n)}(d)$ dla $n = 2, \ldots, 12$
($p_i = i$-ta liczba pierwsza). Kropka ($\cdot$): $K^{(n)}(d) = 0$.
Pozycja $d = 8$: aktywna przy $n=3$, dezaktywowana przy $n=4$.
Pozycja $d = 17$: $K^{(6)}(17) = 2$.
Pozycja $d = 47$: $K^{(12)}(47) = -3$.
Tabela dostarcza empirycznych przykładów wartości
$|K^{(n)}(d)| = 2, 3$, motywujących problem otwarty~(i).
Zjawisko aktywacji i dezaktywacji pozycji,
empirycznie wynikające z kasowania parzystościowego
między nakładającymi się reprezentacjami sum podzbiorów,
jest treścią problemu otwartego~(ii).}
\label{tab:positions}
\end{table}


\section*{A.6. Weryfikacja Hipotezy~\ref{conj:harmonic}}
\label{subsec:a6-hipoteza}

\begin{table}[H]
\centering
\footnotesize
\begin{tabular}{lrrrr}
\toprule
Ciąg generatorów & Pierw.\ gen. & $n^*$
  & $\delta/P$, $n=25$ & $\delta/P$, $n=300$ \\
\midrule
\multicolumn{5}{l}{\textit{Ciągi pierwsze}} \\
Pierwsze $n$ liczb pierwszych &    2 & 8 & 0{,}9998 & 1{,}0000 \\
Liczby pierwsze od $101$      &  101 & 2 & 0{,}9999 & 1{,}0000 \\
Liczby pierwsze od $1009$     & 1009 & 2 & 0{,}9997 & 1{,}0000 \\
Co druga liczba pierwsza ($p_2, p_4, \ldots$) &    3 & 8 & 0{,}9998 & 1{,}0000 \\
Co trzecia liczba pierwsza ($p_3, p_6, \ldots$) &  5 & 2 & 0{,}9998 & 1{,}0000 \\
Mniejsze z par bliźniaczych   &    3 & 9 & 0{,}9996 & 1{,}0000 \\
Liczby pierwsze $\equiv 3 \pmod{4}$ &  3 & 8 & 0{,}9998 & 1{,}0000 \\
Liczby pierwsze $\equiv 1 \pmod{4}$ &  5 & 2 & 0{,}9998 & 1{,}0000 \\
\multicolumn{5}{l}{\textit{Ciągi losowe (losowe liczby nieparzyste)}} \\
Ta sama $H_n$ co ciąg bazowy (seed 42)  & {---} & 8 & 0{,}9999 & {---} \\
Ta sama $H_n$ co ciąg bazowy (seed 137) & {---} & 8 & 0{,}9999 & {---} \\
Bez dopasowania $H_n$ (seed 42)         & {---} & 2 & 0{,}9999 & {---} \\
\bottomrule
\end{tabular}
\caption{Dane są zgodne z Twierdzeniem~\ref{thm:density}:
$\delta(n)/P(n) \to 1$ dla wszystkich testowanych ciągów.
Przy $n=300$: $\delta(n)/P(n) = 1{,}0000$ (sześć miejsc dziesiętnych)
dla wszystkich ciągów pierwszych.
Ciągi losowe o tej samej $H_n$ co ciąg bazowy dają identyczne $n^*$,
co jest zgodne z Hipotezą~\ref{conj:harmonic}.}
\label{tab:sequences}
\end{table}